\documentclass[10pt]{article}
\usepackage{template/template}
\NewsSetKeys{date={2026年11月17日},vol={035},title={每日有趣新闻},author={编辑部},note={待补编短刊：本稿尚未通过图片、构建及PDF验收，不是正式发行版。}}
\begin{document}
\begin{documents}
\newsitem{自主武器的规则，究竟由谁决定？}{国际与科技政策｜11月16日}{联合国《特定常规武器公约》缔约方在日内瓦讨论是否就致命自主武器系统展开谈判。这里的争议并不只是要不要使用人工智能，而是机器能否在缺少逐次人工确认的条件下选择并攻击目标，以及战场上的责任应由谁承担。路透社报道，美国、俄罗斯与以色列反对具有法律约束力的新条约，主张现有人道法框架足够；美国倾向非约束性指引。英国、德国、法国支持讨论法律地位尚待确定的国际文书，会议主席则提出两年内谈判约束性议定书的方案。以上属于各方谈判立场，并非已经生效的国际法。联合国秘书长与红十字会担忧缺少人类控制将增加误伤及追责困难。值得注意的是，武器的自动导航、识别和自主决定致命打击并不是同一技术层级；把所有军用AI混称为自主杀伤武器，反而会模糊真正需要规范的决策边界。谈判是否启动、最后采用条约还是自愿指南，仍未确定。\cite{reuters_laws}}
\newsitem{AI可能如何放大金融系统风险？}{经济与AI｜11月16日}{美联储理事迈克尔·巴尔警告，人工智能应用可能把金融市场既有风险放大为系统性问题。其引用的金融稳定调查中，87\%的受访者将AI列为显著风险；这一比例表示受访者的风险判断，并非87\%的金融机构已经发生AI事故。自动交易策略可能在相似信号下同时买卖，形成羊群效应；在账户代理场景中，银行还需要辨别自动执行指令是否真正获得客户授权。风险并不只来自模型答错题：即使每笔交易在技术上都合法，短时间集中提款也可能使传统的风险缓冲措施来不及反应。巴尔同时提及网络攻击、算法歧视与模型可解释性不足。这是监管者对潜在机制的警示，不能写成已经出现由AI引发的系统性金融危机。\cite{reuters_barr}}
\newsitem{欧洲抢购液化天然气，为什么仍然不够？}{能源经济｜11月16日}{路透社能源评论指出，欧洲临近冬季增加液化天然气进口，但储气仍低于多年同期水平。其给出的11月11日欧盟及英国储气量约810太瓦时，比十年同期均值低约20\%。储气量是能量单位，不能直接与以吨计的月度进口量相加。文章预计11月进口约1011万吨，较10月增长，却低于上年11月。中东冲突与霍尔木兹航运风险限制卡塔尔供应，欧洲采购可能挤压亚洲现货需求。此处属于路透社分析性文章的供需判断，并非国际能源机构发布的确定预测；天气、工业需求和船运恢复都会改变结论。\cite{reuters_lng}}
\newsitem{核裁军会谈尚未等于裁军协议}{外交｜11月16日}{中国、美国和俄罗斯领导人计划在深圳举行会晤。联合国秘书长古特雷斯表示愿意支持可能出现的核裁军努力，同时明确表示结果尚无保证。各方在军控范围、核力量规模及核查机制上仍有分歧。不能把外交议程、谈判意愿或政治表态等同于已经签署的削减核武器协议。\cite{reuters_un}}
\newsitem{航天与地球观测：NASA发布新材料}{科学简讯｜11月16日}{NASA当日科学新闻列出韦布望远镜拍摄的NGC 7129恒星形成区域图像，以及关于印尼查亚峰残存冰体的报道。相关图像值得在取得原始文件并核对署名、拍摄日期与使用条款后收入正式版；本短刊未将网页缩略图冒充已归档的科研原图。\cite{nasa_science}}
\newsitem{数学检索记录}{数学｜11月16日}{本期检索了arXiv数学与Lean相关材料，但未取得可充分核实为11月16日新提交、且已检查主要定理陈述与证明的原始预印本。因此不刊登所谓当日数学突破；待后续补查数论、代数、几何、组合、逻辑、概率与mathlib更新。}
\end{documents}
\small\section*{参考资料}\bibliographystyle{plain}\bibliography{sources/2026/Nov/17/refs}
\end{document}
